\section*{Capítulo \ref{ch:b2:asexual-reproduction} --- Reproducción asexual y clonación en las plantas}

\begin{solution}{exo:b2:asexual-reproduction:1}
\omterm{def:b2:asexual-reproduction:clone}{Clon}: los descendientes genéticamente idénticos de un progenitor
obtenidos por mitosis. \omterm{def:b2:asexual-reproduction:clone}{Geneto}: el individuo genético, todo lo que ha
crecido a partir de un cigoto. \omterm{def:b2:asexual-reproduction:clone}{Rameto}: una unidad fisiológicamente
independiente de un \omterm{def:b2:asexual-reproduction:clone}{geneto}. Órganos: estolones (fresa), rizomas (grama),
tubérculos (patata), bulbos (cebolla), retoños de raíz (álamo temblón),
plántulas foliares (\emph{Kalanchoe}).
\end{solution}

\begin{solution}{exo:b2:asexual-reproduction:2}
Los cámbiums vasculares de la púa y del portainjerto deben estar alineados
y en contacto; el callo tiende un puente sobre el corte y diferencia
xilema y floema a través de él. El portainjerto aporta las raíces (agua,
minerales, vigor, tolerancia al suelo, control del tamaño); la púa, el
brote que fructifica (la variedad). Cada parte conserva su propio genoma
y las células nunca se fusionan: dos \omterm{def:b2:meiosis-heredity:vocabulary}{genotipos} uno junto a otro, una
quimera; un híbrido llevaría una mezcla de ambos genomas en cada célula.
\end{solution}

\begin{solution}{exo:b2:asexual-reproduction:3}
\omterm{def:b2:asexual-reproduction:clone}{Multiplicación vegetativa}: una parte del cuerpo enraíza y da una planta
nueva. \omterm{def:b2:asexual-reproduction:apomixis}{Apomixis}: se forma una \omterm{def:b2:angiosperm-reproduction:seed}{semilla} sin fecundación, cuyo embrión es un
\omterm{def:b2:asexual-reproduction:clone}{clon} de la madre. \omterm{def:b2:asexual-reproduction:apomixis}{Partenogénesis}: un animal se desarrolla a partir de un
óvulo no fecundado. Solo la \omterm{def:b2:asexual-reproduction:apomixis}{apomixis} conserva la \omterm{def:b2:angiosperm-reproduction:seed}{semilla} --- con su
cubierta, su \omterm{def:b2:angiosperm-reproduction:seed}{latencia} y su dispersión.
\end{solution}

\begin{solution}{exo:b2:asexual-reproduction:4}
Se esteriliza el explante; se coloca sobre agar con azúcar, minerales,
vitaminas y auxina y citoquinina equilibradas (callo); se eleva la
proporción citoquinina/auxina para inducir brotes; se subcultivan los
brotes cada pocas semanas; se transfieren a un medio rico en auxina para
que enraícen; se aclimatan las plántulas con humedad elevada y después en
tierra.
\end{solution}

\begin{solution}{exo:b2:asexual-reproduction:5}
Cada planta es sustituida por $1 + 24 = 25$: al cabo de 4 años, $25^{4} =
\num{390625}$ plantas, que necesitan \qty{39000}{m^2}, casi cuatro
hectáreas. El espacio, la luz y los nutrientes, el corto alcance de un
estolón y la muerte de plantas por hacinamiento detienen la fase
exponencial en un par de años; el \omterm{def:b2:asexual-reproduction:clone}{clon} avanza entonces como un frente.
\end{solution}

\begin{solution}{exo:b2:asexual-reproduction:6}
$\num{14000}/130 \approx 110$ generaciones de troncos. Radio de un disco
de \qty{4.3e5}{m^2}: $\sqrt{4.3\times 10^{5}/\pi} = \qty{370}{m}$;
avance de $370/\num{14000} = \qty{0.026}{m}$ al año, menos de tres
centímetros.
\end{solution}

\begin{solution}{exo:b2:asexual-reproduction:7}
$p_t = 2^{t}p_0/(1 + (2^{t} - 1)p_0)$: $t = 5$: $0.0032$; $t = 10$:
$0.1024/1.1023 = 0.093$; $t = 15$: $3.28/4.28 = 0.77$. Nunca se extiende
si su eficacia es inferior a la mitad de la del linaje sexual: una
eficacia relativa $w$ da $q' = 2wq$, que decae para $w <
\tfrac12$.
\end{solution}

\begin{solution}{exo:b2:asexual-reproduction:8}
$10 : 1$: raíces; $1 : 1$: callo indiferenciado; $1 : 10$: brotes.
Un esqueje ya tiene brote y necesita raíces, que la auxina induce en la
base cortada.
\end{solution}

\begin{solution}{exo:b2:asexual-reproduction:9}
1, 10, 100, \num{1000}, \num{10000} kg: tras la tercera temporada solo
hay \qty{1}{t}; tras la cuarta, \qty{10}{t}: cuatro temporadas. La
\omterm{def:b2:angiosperm-reproduction:seed}{semilla} verdadera da miles de \omterm{def:b2:angiosperm-reproduction:seed}{semillas} por planta, pero la patata es un
tetraploide \omterm{def:b2:meiosis-heredity:vocabulary}{heterocigoto}, de modo que las plántulas segregan y ninguna es
la variedad.
\end{solution}

\begin{solution}{exo:b2:asexual-reproduction:10}
$e^{-5} = 0.0067$: 670 individuos libres de \omterm{def:b2:genome-diversification:mutation}{mutaciones} en la población
grande, unos 7 en la pequeña. Siete individuos pueden quedarse todos sin
descendencia por azar en unas pocas generaciones (la probabilidad de que
ninguno se reproduzca es del orden de $e^{-7}$ por generación), tras lo
cual la mejor clase desaparece para siempre y la siguiente pasa a ser la
mejor: un diente del trinquete, que se repite hasta que la carga es
insoportable. Con 670 la pérdida es prácticamente imposible --- hasta que
un cuello de botella reduzca $N$ o aumente la \omterm{thm:b2:genome-diversification:rate}{tasa de mutación}; el
trinquete solo gira en un sentido.
\end{solution}

\begin{solution}{exo:b2:asexual-reproduction:11}
En el \omterm{def:b2:asexual-reproduction:clone}{clon} todos los huéspedes son sensibles: la cepa se propaga sin
freno al ritmo al que viajan las esporas. En la población sexual el
\omterm{def:b2:meiosis-heredity:vocabulary}{genotipo} que vence es una de $2^{10} = 1024$ combinaciones, presente en
alrededor de una milésima parte de las plantas, dispersa entre vecinas
resistentes, de modo que la epidemia se queda sin alimento. La Lumper era
un solo \omterm{def:b2:asexual-reproduction:clone}{clon} cultivado en la mayor parte de un país; el mildiu encontró
todas las plantas indefensas, y no hubo variación con la que mejorar
hasta que se importaron otras variedades.
\end{solution}

\begin{solution}{exo:b2:asexual-reproduction:12}
El \omterm{prop:b2:asexual-reproduction:whysex}{trinquete de Muller}, la combinación de \omterm{def:b2:genome-diversification:mutation}{mutaciones} favorables y la
Reina Roja dan cada uno al sexo un beneficio que puede superar un factor
dos en un mundo grande, parasitado y cambiante. Los rotíferos bdeloideos
son el caso incómodo: cuarenta millones de años sin machos y sin colapso
por el trinquete. Parecen escapar gracias a una resistencia extrema a la
desecación (que mata a sus parásitos), a la captación de ADN ajeno y a
una \omterm{prop:b2:genome-diversification:repair}{reparación del ADN} de una eficacia inusual --- excepciones que
muestran lo que la regla exige normalmente.
\end{solution}

\begin{solution}{pb:b2:asexual-reproduction:1}
\textbf{1.} $4\times 10^{4}$ plantas.
\textbf{2.} $1 + 6\times 2 = 13$.
\textbf{3.} \num{1300}; \num{16900}; \num{219700}.
\textbf{4.} $13^{t} = 400$: $t = \ln 400/\ln 13 = 5.99/2.56 = 2.3$
años --- lleno durante el tercer año.
\textbf{5.} Radio de un disco de una hectárea:
$\sqrt{10^{4}/\pi} = \qty{56}{m}$; a \qty{0.5}{m} al año, \num{113} años.
\textbf{6.} Cien fundadoras dispersas llenan el campo exponencialmente
desde cien frentes en dos años; una sola fundadora llenaría su entorno en
dos años y después avanzaría a paso de tortuga durante un siglo.
\textbf{7.} $5^{n} \ge 10^{6}$: $n \ge 8.6$, nueve subcultivos
($1.95\times 10^{6}$ brotes), 36~semanas.
\textbf{8.} $0.8\times 1.95\times 10^{6} = 1.6\times 10^{6}$ plantas.
\textbf{9.} $0.98$; $0.98^{10} = 0.82$.
\textbf{10.} $10\times 0.98 = 9.8$ líneas limpias esperadas; de diez
esquejes, una.
\textbf{11.} El virus avanza por los plasmodesmos y el floema, y la cúpula
apical no tiene conexión vascular y se divide más deprisa de lo que se
propaga el virus, de modo que la décima de milímetro más externa suele
estar libre de infección.
\textbf{12.} Laboratorio: plantas uniformes y sin virus de inmediato, pero
costosas, y cualquier \omterm{def:b2:genome-diversification:mutation}{mutación} o contaminación del cultivo se multiplica
un millón de veces. Estolones: gratuitos y sobre el terreno, pero se
pierden tres temporadas y se arrastran todos los virus de las plantas
madre.
\textbf{13.} Los apomícticos son $pN$ y dan $F$ \omterm{def:b2:angiosperm-reproduction:seed}{semillas} cada uno; los
sexuales, $(1 - p)N$, dan $F/2$ \omterm{def:b2:angiosperm-reproduction:seed}{semillas} cada uno (la mitad de su
esfuerzo es \omterm{def:b2:plant-life-cycles:heterospory}{polen}, que no produce \omterm{def:b2:angiosperm-reproduction:seed}{semillas}); la parte de \omterm{def:b2:angiosperm-reproduction:seed}{semillas} de los
apomícticos es $Fp/(Fp + F(1 - p)/2)
= 2p/(1 + p)$.
\textbf{14.} $p_5 = 0.032/1.031 = 0.031$; $p_{10} = 1.024/2.023 =
0.51$: supera la mitad en la generación 10.
\textbf{15.} $p' = 2p(1 - cp)/\bigl(2p(1 - cp) + (1 - p)\bigr)$.
\textbf{16.} $p' = p$ cuando $2(1 - cp^*) = 1$: $p^* = 1/2c = 0.625$.
Por debajo de $p^*$, $2(1 - cp) > 1$ y $p$ aumenta; por encima,
disminuye: estable.
\textbf{17.} $p^* \ge 1$ cuando $c \le \tfrac12$: el apomíctico se impone
salvo que el parásito le cueste, cuando es común, más de la mitad de su
producción.
\textbf{18.} $c = 0.6$: $p^* = 0.83$. Para mantener el sexo en la
población, el parásito debe, cuando el apomíctico es frecuente, reducir la
eficacia del \omterm{def:b2:asexual-reproduction:clone}{clon} más que la doble ventaja --- la Reina Roja tiene que
correr deprisa.
\textbf{19.} Tres juegos de cromosomas no pueden aparearse en la \omterm{def:b2:meiosis-heredity:meiosis}{meiosis},
de modo que sus gametos están desequilibrados; queda excluido de la
recombinación, y el \omterm{prop:b2:asexual-reproduction:whysex}{trinquete de Muller} y los parásitos actúan sobre él
sin oposición --- por eso los linajes apomícticos de diente de león son
jóvenes y se forman una y otra vez a partir de antepasados sexuales.
\textbf{20.} $10^{-9}\times 4.8\times 10^{8} = 0.48$ por división.
\textbf{21.} $2\times \num{14000}\times 20 = 5.6\times 10^{5}$
divisiones; $0.48\times 5.6\times 10^{5} = 2.7\times 10^{5}$
\omterm{def:b2:genome-diversification:mutation}{mutaciones}.
\textbf{22.} $2.7\times 10^{5}/4.8\times 10^{8} = 5.6\times 10^{-4}$ del
genoma; unas \num{5400} diferencias codificantes.
\textbf{23.} La competencia entre linajes celulares en el meristemo (una
línea mutante que se divide más despacio es desplazada) y entre \omterm{def:b2:asexual-reproduction:clone}{rametos}
(un tronco portador de una \omterm{def:b2:genome-diversification:mutation}{mutación} perjudicial emite menos retoños). El
filtro es más débil porque la mayoría de las \omterm{def:b2:genome-diversification:mutation}{mutaciones} somáticas son
heterocigotas y quedan enmascaradas, y nada las expone: ninguna \omterm{def:b2:meiosis-heredity:meiosis}{meiosis}
las vuelve homocigotas y ninguna fase de cigoto \omterm{def:b2:unicellular-diversity:unicellular}{unicelular} pone a prueba
el genoma entero en una sola célula.
\textbf{24.} No es uniforme: sus troncos difieren en cientos de miles de
sitios. Pero las diferencias son en su mayoría neutras y están
enmascaradas, de modo que el \omterm{def:b2:asexual-reproduction:clone}{clon} es en la práctica un solo \omterm{def:b2:meiosis-heredity:vocabulary}{genotipo} ---
un mosaico de variantes somáticas de un mismo \omterm{def:b2:asexual-reproduction:clone}{geneto}.
\textbf{25.} Campo lleno al cabo de 2,3 años; un meristemo da
$1.6\times 10^{6}$ plantas en nueve meses; con $c = 0.8$ el apomíctico
se estabiliza en $p^* = 0.625$, y los sexuales conservan el
\qty{37.5}{\%}.
\end{solution}
